%Usage guide: https://docs.google.com/document/d/1oDJKcFVK_yXD3Do6vfyTFKJuDwdwN4g6XFTM5QASadA/edit?usp=sharing

\documentclass[12pt]{article} % If you want 12 pt font instead of the standard 10 pt
%\documentclass{article}
\usepackage[utf8]{inputenc}


\usepackage{amsmath}
\usepackage{amssymb}
\usepackage{stackengine}





\usepackage[shortlabels]{enumitem}
\usepackage[margin=1 in]{geometry}

\newcommand{\LRT}[2]{%
  \mathrel{\mathop\gtrless\limits^{#1}_{#2}}%
}


\pagenumbering{arabic}

\usepackage{mathhw}

\begin{document}

\setlength{\abovedisplayskip}{7pt}
% \setlength{\belowdisplayskip}{10pt}

\begin{center}
    \textbf{6.437 Notes}\\
    Anugrah Chemparathy
\end{center}
\note{L9 - Exponential Families}

\begin{align*}
    \dot{\alpha}(x) 
    &= \dot{\lambda}(x) E[t(y)] \\
    \Ddot{\alpha}(x)
    &= \Ddot{\lambda}(x) E[t(y)] + \dot{\lambda}(x)^2\ \Var[t(y)]
\end{align*}

\begin{equation*}
    J_y(x) =
    E\left[\left(\frac{\partial}{\partial x} \ln p_y(y;x) \right) ^2 \right]
    =- E\left[\frac{\partial^2}{\partial x^2} \ln p_y(y;x) \right]
\end{equation*}

\note{L18 - Conjugate Priors}
\begin{equation*}
    p_{x|y_N,\cdots,y_1} = 
    \frac{
    p_{y_N|x}p_{x|y_{N-1}, \cdots, y_1}
    }{
    \int p_{y_N|x} p_{x|y_{N-1}, \cdots, y_1} \ dx
    }
\end{equation*}

In one of the optional sections, it is shown that conjugate priors only exist for $p_{y|x}$ which are from exponential families. In the last section, they show how to construct a conjugate prior family from a given exponential family, and essentially that all exponential families have a conjugate prior family.

If $p_{y|x}(y|x) = \exp{\lambda(x)^T t(y) - \alpha(x) + \beta(y)}$, then the corresponding conjugate prior family is:
\begin{equation*}
    Q = \{ 
    q(\cdot; \mathbf{t},N):
    q(x; \mathbf{t},N)
    =
    \exp{[t^T \lambda(x) - N\alpha(x)] - \gamma(t,N)}
    \}
\end{equation*}
where $\mathbf{t}$ is an updating function of the $t(y)$ we see, starting from our prior's value of $t_0$. The corresponding prior is $p_x(x) = q(x;t_0,N_0)$.

Correspondingly, $Q$ is itself a canonical exponential family with natural statistic $[\lambda(x), - \alpha(x)]$ and parameter (known as $x$ in the definition earlier in the notes) $[\mathbf{t}, N]$

\begin{align*}
    p_{y|x}(y|x) &= \exp{\lambda(x)^T t(y) - \alpha(x) + \beta(y)} \\
    Q &= \{ 
    q(\cdot; \mathbf{t},N):
    q(x; \mathbf{t},N)
    =
    \exp{[t^T \lambda(x) - N\alpha(x)] - \gamma(t,N)}
    \}
\end{align*}

\note{general}
\begin{equation*}
    P(\cup_N A_i) \leq \sum_N P(A_i)
\end{equation*}

\end{document}


\begin{comment}
Pink - Important ideas/defs/thsm
Green - Corollaries/lesser thms 
Blue - Smaller words
Yellow - Examples
Gray - Anything trivial

\begin{itemize}[topsep=-1pt, itemsep=-1pt]
    \item 
\end{itemize}


\begin{enumerate}[label=\textbf{(\alph*)}, topsep=0pt]
    \item Hello
\end{enumerate}

\begin{itemize}[topsep=0pt]
    \item Hi
\end{itemize}

\begin{center}
   \includegraphics[width=0.6\textwidth]{L15-1.png} 
\end{center}
\end{comment} 